\documentclass[output=paper,colorlinks,citecolor=brown]{langscibook}
\ChapterDOI{10.5281/zenodo.20541864}
\author{Kultida Khammee\orcid{}\affiliation{University of Phayao}}
\title{Grammaticalization of future markers in Thai}
\abstract{Contemporary Thai has a number of forms that mark the notion of futurity. The primary future tense marker is \textit{cà}, whose lexical origin, though not yet established, is likely the verb \textit{càk} ‘know’. Other markers are used either alone or in combination with \textit{cà(k)}\textit{.} Their morphosyntactic development reveals some influence of the typological and idiosyncratic features of Thai, such as pragmatic orientation, verb serialization, and preference for polylexemic units in lexicalization and grammaticalization, among others. Their semantic developments are also well motivated in terms of conceptual relatedness between the source forms and the grammaticalized forms. A review of grammaticalization scenarios with reference to grammaticalization parameters, e.g., desemanticization, extension, decategorialization, and erosion, shows that the Thai future markers have been grammaticalized to varying degrees.}

%move the following commands to the "local..." files of the master project when integrating this chapter


\IfFileExists{../localcommands.tex}{
   \addbibresource{../localbibliography.bib}
   \usepackage{tabularx,multicol}
\usepackage{url}
\urlstyle{same}

\usepackage[english]{babel}
 
\usepackage{langsci-optional}
\usepackage{langsci-lgr}
\usepackage{langsci-gb4e}

\usepackage{langsci-branding}

%\usepackage{tabularx}

% elena
\usepackage{graphicx} %package to manage images
\usepackage{paralist}

% sarah & lena
\usepackage[justification=centering]{caption}
\usepackage{rotating}

% robin
\usepackage{lscape}
\usepackage{tikz}
\usetikzlibrary{tikzmark}

% dominika
\usepackage{multirow}
\usepackage{hanging}

% olaf & dylan
\usepackage{subcaption}

% svetlana
\usepackage{verbatim}
\usepackage{enumitem}

% wiemer
%\usepackage[markup=underlined]{changes}
%\usepackage{hyphenat}

% wir
\usepackage{longtable}

% for crossref in the introduction jiabin
\usepackage{xr}
   % \newcommand*{\orcid}{}
%

\makeatletter
\let\thetitle\@title
\let\theauthor\@author
\makeatother

\newcommand{\togglepaper}[1][0]{
   \bibliography{../localbibliography}
   \papernote{\scriptsize\normalfont
     \theauthor.
     \titleTemp.
     To appear in:
     E. Di Tor \& Herr Rausgeberin (ed.).
     Booktitle in localcommands.tex.
     Berlin: Language Science Press. [preliminary page numbering]
   }
   \pagenumbering{roman}
   \setcounter{chapter}{#1}
   \addtocounter{chapter}{-1}
}
%
\newbool{bookcompile}
\booltrue{bookcompile}
\newcommand{\bookorchapter}[2]{\ifbool{bookcompile}{#1}{#2}}
%
% %\newcommand{\orcid}[1]{}
\graphicspath{ {./figures/} }
%
\let\eachwordone=\itshape


% Seongha
% \newcommand{\koreex}[1]{{\fontspec{Baekmuk Batang}{#1}}}

% Kutida
\newfontfamily\thaifont{Noto Sans Thai} 
\newcommand{\textthai}[1]{{\thaifont #1}} % für Thai


\newfontfamily\chinesefont{Noto Sans CJK SC}
\newcommand{\textchinese}[1]{{\chinesefont #1}} % Für Chinesisch

% wiemer
\def\hyph{-\penalty0\hskip0pt\relax}
% \definechangesauthor[color=NavyBlue]{SH}



\newcolumntype{K}{>{\raggedright\arraybackslash}p{0.6\textwidth}}  % Right column with 0.7 width


% \newfontfamily\krn[Scale=MatchLowercase,BoldFont=SourceHanSerif-Bold.otf]{SourceHanSerif-Regular.otf}
% \AdditionalFontImprint{Source Han Serif KO}
% \XeTeXlinebreaklocale 'ko'
% \renewcommand{\krn}{}



% Rhee
\newfontfamily\koreanfont{Noto Sans CJK KR} 
\newcommand{\textkorean}[1]{{\koreanfont #1}} % für Koreanisch

\renewcommand{\REF}[2][]{(\ref{#2}#1)}
   %% hyphenation points for line breaks
%% Normally, automatic hyphenation in LaTeX is very good
%% If a word is mis-hyphenated, add it to this file
%%
%% add information to TeX file before \begin{document} with:
%% %% hyphenation points for line breaks
%% Normally, automatic hyphenation in LaTeX is very good
%% If a word is mis-hyphenated, add it to this file
%%
%% add information to TeX file before \begin{document} with:
%% %% hyphenation points for line breaks
%% Normally, automatic hyphenation in LaTeX is very good
%% If a word is mis-hyphenated, add it to this file
%%
%% add information to TeX file before \begin{document} with:
%% \include{localhyphenation}

\hyphenation{
    con-tri-bute
    par-a-digm
    pe-ri-phra-ses
    achieve-ment
    accomp-lish-ment
    We-li-nchen-kang-ci-kok
    cross-ling-u-is-tic
    rec-ord-s
    trun-ca-ted
    de-ri-va-tion-al
    stem-de-ri-va-tion-al
    pre-sent
    in-flu-en-ced
    sys-tem
    ques-tions
    Bor-kovs-kij
    Da-ny-len-ko
    e-qui-va-lents
    re-cog-nized
    In-ter-fe-renz-schicht
    Aus-ei-nan-der-se-tz-ung
    Früh-neu-hoch-deutsch-kor-pus
    Ak-tions-art
    Smir-no-va
    Whi-le
    Ar-chive
    bi-ginnan
    Rei-chen-bach
    ger-ma-nis-ti-sche
    Re-fe-renz-kor-pus
    le-xi-ka-lisch-en
    Alt-hoch-deutsche
    Li-te-ra-tur
    dia-chron
    Dia-lek-to-lo-gie
}

\hyphenation{
    con-tri-bute
    par-a-digm
    pe-ri-phra-ses
    achieve-ment
    accomp-lish-ment
    We-li-nchen-kang-ci-kok
    cross-ling-u-is-tic
    rec-ord-s
    trun-ca-ted
    de-ri-va-tion-al
    stem-de-ri-va-tion-al
    pre-sent
    in-flu-en-ced
    sys-tem
    ques-tions
    Bor-kovs-kij
    Da-ny-len-ko
    e-qui-va-lents
    re-cog-nized
    In-ter-fe-renz-schicht
    Aus-ei-nan-der-se-tz-ung
    Früh-neu-hoch-deutsch-kor-pus
    Ak-tions-art
    Smir-no-va
    Whi-le
    Ar-chive
    bi-ginnan
    Rei-chen-bach
    ger-ma-nis-ti-sche
    Re-fe-renz-kor-pus
    le-xi-ka-lisch-en
    Alt-hoch-deutsche
    Li-te-ra-tur
    dia-chron
    Dia-lek-to-lo-gie
}

\hyphenation{
    con-tri-bute
    par-a-digm
    pe-ri-phra-ses
    achieve-ment
    accomp-lish-ment
    We-li-nchen-kang-ci-kok
    cross-ling-u-is-tic
    rec-ord-s
    trun-ca-ted
    de-ri-va-tion-al
    stem-de-ri-va-tion-al
    pre-sent
    in-flu-en-ced
    sys-tem
    ques-tions
    Bor-kovs-kij
    Da-ny-len-ko
    e-qui-va-lents
    re-cog-nized
    In-ter-fe-renz-schicht
    Aus-ei-nan-der-se-tz-ung
    Früh-neu-hoch-deutsch-kor-pus
    Ak-tions-art
    Smir-no-va
    Whi-le
    Ar-chive
    bi-ginnan
    Rei-chen-bach
    ger-ma-nis-ti-sche
    Re-fe-renz-kor-pus
    le-xi-ka-lisch-en
    Alt-hoch-deutsche
    Li-te-ra-tur
    dia-chron
    Dia-lek-to-lo-gie
}
   \boolfalse{bookcompile}
   \togglepaper[11]%%chapternumber
}{}

\begin{document}
\maketitle


\section{Introduction}

Living in the spatio-temporal dimension, all humans supposedly have linguistic means, be they lexical or grammatical, of expressing future time reference.\footnote{For instance, \citet[108--109]{Dahl1985} notes that even in the languages without a future tense marker, there are ways of signaling future time reference.} As observed in the literature, Thai does not have an inflectional future, as a general characteristic of an isolating language, but rather has a construction involving the auxiliary \textit{cà} combining with a verb to indicate future time reference \citep[173--174]{Dahl1985}.\footnote{Thai linguists use a few different romanization systems. We use the Thai IPA transcription.} Thai also has a handful of other markers, with variable degrees of grammaticalization, that signal that an event or state of affairs will occur or exist at a future time. Their developmental patterns exhibit certain commonalities with other languages as well as peculiarities due to the typological or idiosyncratic features of Thai. This chapter describes the means of future time reference in Thai from a grammaticalization perspective. The data sources include the Thai National Corpus (TNC), dictionaries, lexica, online resources, and reference grammars.\footnote{The Thai National Corpus (TNC) is an online searchable, 33.4-million-word corpus, developed by Chulalongkorn University.}

Thai is a Kra-Dai language (\cite{Ostapirat2000}; \cite{Eberhard2022}), also known as Tai-Kadai (\cite{DillerEdmondsonLuo2008}), and is mostly spoken by the people in Thailand and its vicinity. Thai is an isolating language. Word order is SVO and phrases are head-initial. Thai does not have case marking morphology, but it has a complex tone system (for typological descriptions, see \cite{Noss1954}, \citeyear{Noss1964}; \cite{abramson1962}; \cite{Haas1964}; \cite{anthony1968}; \cite{Warotamasikkhadit1972}; \cite{TingsabadhAbramson1999}; \cite{Campbell2000}; among others). 

% \largerpage[1.5]

This chapter is organized as follows: Section \ref{sec:khammee:2} describes and exemplifies a few forms indicating future time reference, i.e., \textit{càk} (Section \ref{sec:khammee:2.1}), \textit{cà} (Section \ref{sec:khammee:2.2}), adverbials (Section \ref{sec:khammee:2.3}), and particles (Section \ref{sec:khammee:2.4}); Section \ref{sec:khammee:3} addresses some prominent issues with respect to grammaticalization of future time reference, focusing on the morphosyntactic motivation (Section \ref{sec:khammee:3.1}), the conceptual motivation (Section \ref{sec:khammee:3.2}), and grammaticalization parameters (Section \ref{sec:khammee:3.3}); and Section \ref{sec:khammee:4} summarizes the findings and concludes the chapter. 


\section{Future time reference in Thai} \label{sec:khammee:2}

Some Thai reference grammars explicitly state that Thai has no future tense markers, largely because their use is not mandatory and future time reference can be inferred from a temporal adverb, an aspect marker, or from the context (\cite{IwasakiIngkaphirom2005}: 123; \cite{Thiengburanathum2013}: 159). However, most reference grammars acknowledge the presence of future markers (\cite{Noss1964}: 165; \cite{CampbellShaweevongse1972}: 26; \cite{Setthapun1992}: 24; \cite{Smyth2002}: 67; among others). The most commonly cited future marker is \textit{cà}, but there are other forms as well. The notion of futurity signaled by certain words is often interpreted as a lexical meaning rather than a grammatical function. This is the case because of categorial fluidity, resistance to formal change, as well as the prevalence of the polyfunctionality of forms across lexical and grammatical domains, which are characteristics typical of isolating languages.

However, if we adopt a dynamic perspective such as the one in grammaticalization theory, Thai indeed has a number markers of futurity. They include \textit{càk}\textit{,} \textit{cà}\textit{,} \textit{cuan,} \textit{kʉ̀ap, thɛ̂ɛp, }and \textit{klây}. There are poly-lexemic forms created by means of compounding. Some authors, e.g., \citet{Supanvanich1973}, also claim that certain sentence ending particles in Thai function as markers of future time reference. In the following, we illustrate the uses of such future markers. 

\subsection{\textit{Càk}}\label{sec:khammee:2.1}

% \largerpage[1.5]

The word \textit{càk} marks the future in Modern Thai. Evidently, it is the earlier form of \textit{cà }(see Section \ref{sec:khammee:3.2} below), the primary future marker in Contemporary Thai. Following the tenets of grammaticalization theory according to which all grammatical forms originate from a lexical source, \textit{càk }is also believed to had have a lexical origin. With a regrettably shallow history of documentation (see Section \ref{sec:khammee:3.3}), however, there has not been any serious scholarly investigation to date into its origin, and, consequently, there is no consensus regarding its lexical origin. Despite the paucity of lexical evidence, some clues do exist in language contact. For instance, the Thai lexicon was greatly influenced by Middle Chinese, and a lexicon of Middle Chinese borrowings (\cite{SEALANG}) includes Thai \textthai{จัก} \textit{càk} (\textit{jàk} in its notation). As the historical source of the Thai \textit{càk}, this lexicon points to the Middle Chinese lexeme denoting ‘recognize (as is)’, which references Proto-Southwestern Tai \textit{*cak} ‘to recognize, know’ (\cite{Jonsson1991}), as well as Proto-Southwestern Tai \textit{*čak} and Proto-Tai \textit{*čak }(\cite{Li1977}). The same source lists the Thai verb \textthai{รู้จัก} \textit{rúucàk} (=\textit{ru}\textit{jàk}) ‘to recognize, know’ as a borrowing of the Middle Chinese lexeme \textchinese{識} (\textit{śik}) ‘to recognize, know, discern, have knowledge, etc.’, citing \citeauthor{manomaivibool1975} (\citeyear{manomaivibool1975}; known as ‘Prapin’), \citet{baxter2000}, and \citet{Starostin} (\cite{SEALANG}). Based on these comparative-linguistic analyses, it can be hypothesized that the future marker \textit{càk} (and its later form \textit{cà, }see Section \ref{sec:khammee:2.2} below) developed from the Middle Chinese borrowing \textit{càk} with the meaning ‘know’.\footnote{Other possibilities from the data in \citeauthor{SEALANG} are the two homophonous lexemes (\textthai{จัก} \textit{càk}; \textit{jàk} in its notation), one denoting ‘weave, knit’ (Middle Chinese \textit{ćik/tsyik}) and the other denoting ‘split, pare bamboo’ (Proto Southwestern \textit{*cak/*čak} and Proto-Tai *\textit{čak}), the latter being proposed in the earlier version of this paper. An anonymous reviewer raised concerns that the proposal is not substantiated by historical data. Dr. Weera Ostapirat (p.c.), a specialist in the Tai language family, also expresses his opinion that ‘split, pare bamboo’ is not a likely source and that ‘know’ is the better candidate (see Section \ref{sec:khammee:3.2} for conceptual motivation across languages). He further states that the native Thai verb \textit{càak }‘depart, leave’ is a less likely but possible candidate. At the current level of understanding, the lexical source of the future marker \textit{cà(k)} remains inconclusive. Special thanks go to the reviewer and Dr. Weera Ostapirat for their suggestions.} Despite \textit{càk} being used as the cognitive verb ‘to know, understand, be aware of’, it has become obsolete and only survives in the form of \textit{rúucàk} ‘to know a person, place, etc.’ (thus, recognition-oriented), in which \textit{rúu} also means ‘to know a subject or piece of learned information’ (\href{thai-language.com}{thai-language.com}) (thus, comprehension-oriented).\footnote{It has been observed that Thai has an idiosyncratic tendency of favoring a morphological process of ‘recycling’, whereby synonyms or near synonyms are compounded. The multiple forms thus created tend to create ‘layering’ (\cite{Hopper1991}) and have different degrees of intensification or different register-specific, semantic specialization (see below for more instances). It is also common that there are no noticeable semantic differences between the multiple forms created through recycling. Some of such (near-)synonymous forms may disappear, as shown here with \textit{càk} ‘to know’.} The use of the reinforced \textit{rúucàk }is exemplified in~\xref{ex:1}:\footnote{Note that Thai does not practice inter-lexical spacing, but Thai data are given with spacing for visual clarity.}


\ea \label{ex:1} Modern Thai (constructed)\\
    \textthai{ผู้ชาย} \textthai{ที่} \textthai{ผึ้ง} \textthai{รู้จัก} \textthai{ใน} \textthai{เมือง} \textthai{เชียงใหม่} \textthai{ก็} \textthai{มี} \textthai{อยู่} \textthai{สอง} \textthai{คน}\\
         \gll \textit{phûuchaay} \textit{thîi} \textit{phʉ̂ŋ} \textit{rúucàk} \textit{nay} \textit{mʉaŋ} \textit{chiaŋ.mày} \textit{kɔ̂ɔ} \textit{mii} \textit{yùu} \textit{sɔ̌ɔŋ} \textit{khon}\\
         man that [name] know in city [name] also have be two person\\
    \glt ‘There are two men whom Pheung knows in Chiang Mai.’ \\
    \hfill (Modern Thai, constructed)
\z

\textit{Càk} has been grammaticalized as a future auxiliary (for morphosyntactic and conceptual motivation see Section \ref{sec:khammee:3.1} and \ref{sec:khammee:3.2}). As an auxiliary, first attested in the Sukhothai period (1238--1347 CE), \textit{càk }indicates future or intention.\footnote{In the excerpt, the parts that the original Thai data do not need to be presented with interlinear morphemic glossing are presented in English translation within square brackets in the interest of space. In the English translation, the parts that are required for syntactic and semantic reasons but are only implicit in the original data are presented within parentheses. An alternative translation is presented in parentheses as well. Also for convenience, markers of future time, regardless of their grammatical status or additional semantics, are glossed as \textsc{fut}(ure), unless a different label is needed for expositional clarity.} Its use is exemplified below:

\largerpage[-1]
\newpage

\ea \label{ex:2} \textit{Càk} future/intention\\
\ea \label{ex:2a}
    \textthai{ใคร} \textthai{จัก} \textthai{ใคร่} \textthai{ค้า} \textthai{ช้าง} \textthai{ค้า} \textthai{ใคร} \textthai{จัก} \textthai{ใคร่} \textthai{ค้า} \textthai{ม้า} \textthai{ค้า} \textthai{ใคร} \textthai{จัก} \textthai{ใคร่} \textthai{ค้า} \textthai{เงือน} \textthai{ค้า} \textthai{ทอง} \textthai{ค้า}\\
         \gll \textit{khray} \textit{càk} \textit{khrây} \textit{kháa} \textit{cháaŋ} \textit{kháa} \textit{khray} \textit{càk} \textit{khrây} \textit{kháa} \textit{máa} \textit{kháa} \textit{khray} \textit{càk} \textit{khrây} \textit{kháa} \textit{ŋʉan} \textit{kháa} \textit{thɔɔŋ} \textit{kháa}\\
         who \textsc{fut} want sell elephant sell who \textsc{fut} want sell horse sell who \textsc{fut} want sell silver sell gold sell\\
    \glt ‘Whoever would want to trade in elephants does so. Whoever would want to trade in horses does so. Whoever would want to trade in silver or gold does so.’ \\ \hfill (1292, King Ramkhamhaeng’s free trade law)
    
    \ex \label{ex:2b} \textit{Càk} future/intention\\
    \textthai{เมื่อ} \textthai{ไม่} \textthai{รัก} \textthai{จัก} \textthai{หย่า} \textthai{ก็} \textthai{ว่า} \textthai{เถิด}\\
         \gll \textit{mʉ̂a} \textit{mây} \textit{rák} \textit{càk} \textit{yàa} \textit{kɔ̂ɔ} \textit{wâa} \textit{thə̀ət}\\
         when not love \textsc{fut} divorce so say please\\
    \glt ‘When (you) do not love and (you) would (intend to) divorce (me), please tell (me).’\\
    \hfill (2009, Enak Phra Khanong)
\z
\z

Excerpt \xref{ex:2a} is taken from a historical text, in which\textit{ càk} functions as a future marker with the main verb \textit{khrây} ‘want’. While \textit{càk} can mean either future or intention (the notion of intention may have arisen from the knowledge meanings associated with the source verb), the presence of the co-occurring conative verb \textit{khrây} ‘want’ makes the intention interpretation redundant and the future interpretation contextually favored. Excerpt \xref{ex:2b} is taken from a Modern Thai source, part of a traditional Thai octameter poem entitled ‘The Legend of Mae Nak Phra Khanong’, which still bears many characteristic features of the classical poem. In this excerpt, \textit{càk} can mean either future or intention. One might argue that if the phrase \textit{càk yàa }‘will divorce/intend to divorce’ is viewed in syntactic parallelism with the preceding \textit{mây rák }‘not love’, \textit{càk} would favor the intention interpretation, i.e., as a main verb, rather than the future interpretation, i.e., as an auxiliary. However, since Thai does not require syntactic parallelism, there is no syntactically-driven interpretational favoring. The future interpretation is as much favored as in \xref{ex:2a}, especially because the second predicate (i.e., ‘divorce’) is later in time (thus, future) relative to the first predicate (i.e., ‘not love’). 

The future auxiliary \textit{càk} can form a deontic modal auxiliary with \textit{tɔ̂ŋ} in the form of \textit{càk-tɔ̂ŋ }‘must’, especially in formal registers, e.g., legal documents. As briefly alluded to before, ‘recycling’ a formant for use in combination with another formant is a common lexicalization-grammaticalization strategy in Thai (\cite{RheeKhammee2022}, \cite{KhammeeRhee2022}, \cite{ThirakunkovitRhee2023}; see also  Section \ref{sec:khammee:3.3.4}). The participating form \textit{tɔ̂ŋ} is itself the deontic modal ‘must’, which is also heterosemous (\cite{Lichtenberk1991}) with the verb \textit{tɔ̂ŋ} ‘to make contact, be exposed to, touch, agree with, etc.’ 

\ea \label{ex:3} càk-tɔ̂ŋ future/deontic necessity\\
    \textthai{ความ} \textthai{ยุติธรรม} \textthai{จักต้อง} \textthai{ดำรง} \textthai{อยู่} \textthai{ใน} \textthai{มวล} \textthai{มนุษย์}\\
         \gll \textit{khwaam} \textit{yútìtham} \textit{càk-tɔ̂ŋ} \textit{damroŋ} \textit{yùu} \textit{nay} \textit{muan} \textit{manút}\\
         \textsc{nmlz} justice \textsc{fut} maintain be in entire human\\
    \glt ‘Justice will (shall, must) prevail in mankind.’\\
    \hfill (Modern Thai, constructed)
\z


Excerpt \xref{ex:3} contains \textit{càk-tɔ̂ŋ}, a future marker also with a deontic necessity meaning. An interesting aspect is that, due to the fact that deonticity is inherently future-related, the presence of the future marker \textit{càk} is not semantically consequential, i.e., there is no perceivable difference between \textit{càk-tɔ̂ŋ} and \textit{tɔ̂ŋ}, both denoting ‘must’. The close relationship between deontic modality (agent-oriented) and future tense has been observed in \citet[]{Fleischman1982}, \citet[]{Dahl1985}, \citet[]{Heine1993}, \citet[]{BybeePerkinsPagliuca1994}, among others. 


\subsection{\textit{Cà}} \label{sec:khammee:2.2}

% \largerpage[-2]

In Modern Thai, the primary future marker is \textit{cà}, which is historically an eroded form of \textit{càk}, addressed in the previous section.\footnote{The older \textit{càk} is now used with very low frequency (52 pmw in a contemporary corpus TNC) and functions as an auxiliary, marking the future tense or the subject's intention. On the other hand, the newer and eroded \textit{cà} occurs with very high frequency as a general future-tense marker (15,014 pmw in TNC).} The future-marking function of \textit{cà} is first identified in the Ayutthaya Period (1351--1767). In the earliest record of \textit{cà}, the form strongly suggests intention, thus often making the sentence ambiguous as to futurity or intention (note that the two are conceptually close, as exemplified in the development of the English future \textit{will}, \cite{Aijmer1985}: 11--20).\footnote{\citet{Soithurum2010}, conceding that the absolute lexical origin of \textit{cà} is unavailable, asserts that its earliest recognizable meaning is ‘intend’.} As indicated previously, it is quite possible that \textit{cà}, from its earlier verbal meanings of \textit{càk} ‘to know’, has developed into ‘to intend’, and its conceptually close ‘future’ meaning has emerged simultaneously or soon afterwards. Thus, even from the early usage, \textit{cà} functions as a marker not only of the future but also of intention and even desire, and other meanings such as assurance and obligation become associated with \textit{cà}, as these notions are all closely related to futurity (cf. \cite{Dahl1985}: 103; also see below). The following is an excerpt involving the future marker \textit{cà}:

\ea \label{ex:4} \textit{cà} future/intention\\
(in a letter by the Siamese ambassador to France)\\
    \textthai{จะ} \textthai{ขอ} \textthai{แก่} \textthai{ท่าน} \textthai{ให้} \textthai{เข้า} \textthai{ทูล} \textthai{แก่} \textthai{พระมหากษัตรเจ้า} \textthai{ผู้} \textthai{ใหญ่} \textellipsis\\
         \gll \textit{cà} \textit{khɔ̌ɔ} \textit{kɛ̀ɛ} \textit{thân} \textit{hây} \textit{khâw} \textit{thuun} \textit{kɛ̀ɛ} \textit{phrámáhǎakàsàt-câaw} \textit{phûu} \textit{yày} [...] \textit{lɛɛ} \textit{cà} \textit{ù-sàa} \textit{thamhây} \textit{cháná} \textit{thúkkhon} \textit{nay} \textit{khwaam} \textit{nóp-nɔ̂ɔp} \textit{nán}\\
         \textsc{fut} beg to you:\textsc{hon} to enter tell to king who great {} and \textsc{fut} attempt make win everybody in \textsc{nmlz} veneration that\\
    \glt ‘We would (intend to) beg Your Excellency to tell His Majesty [...] and (we) will (intend to) endeavor to surpass everybody else in terms of respect and veneration toward the king.’\\
    \hfill (Siamese Documents of the 17\textsuperscript{th} Century, Paris)
\z

Excerpt \xref{ex:4} is taken from a letter dated from the 17\textsuperscript{th} century, in which the Siamese ambassador requests the addressee to convey a message of allegiance and veneration to the French king. As seen in the excerpt, the future marker \textit{cà}, which, in addition to being a future auxiliary, has potential as a verb, and is followed by three successive verbs \textit{ù-sàa} ‘attempt’, \textit{thamhây} ‘make’, and \textit{cháná} ‘win’.\footnote{This is a unique situation in Thai, where multiple verbs may occur in succession asyndetically. For instance, \citet[160]{Diller2006} notes that Thai, which he called ‘a verb-loving language’, allows such series of verbs as ‘Peel give eat’ (for ‘Peel the oranges for the guests’), ‘Tell finish not have’ (for ‘I already told you, there aren’t any’), and the like. See 3.1 for more discussion on its relevance to grammaticalization.}

The auxiliary \textit{cà} denotes, in addition to future tense, various future-related modal notions, e.g., intention, promise, necessity, possibility, probability, certainty, indeterminacy, proximation, constancy, etc., either alone or through ‘univerbation’ (\cite{Lehmann2015}) with other lexemes.\footnote{Since modal markers in Thai are many and varied, with tense and aspect markers omissible without affecting grammaticality, \citet[114]{Soithurum2010} asserts that Thai is not a tense- or aspect-prominent language but a ‘mood-prominent’ language.} These notions are so closely related that whichever is denoted by this particular instance of a \textit{cà}-sentence largely depends on the context, and functional designation may not be exclusive. The usages of \textit{cà} for denoting these modal notions are exemplified below:

\ea \label{ex:5} 
    \ea \label{ex:5a} \textit{cà} future/intention\\
        \textthai{ฉัน} \textthai{จะ} \textthai{เข้าร่วม} \textthai{ประชุม} \textthai{ที่} \textthai{เยอรมนี} \textthai{เดือน} \textthai{หน้า}\\
        \gll \textit{chǎn} \textit{cà} \textit{khâwrûam} \textit{pràchum} \textit{thîi} \textit{yəəramánii} \textit{dʉan} \textit{nâa}\\
        I \textsc{fut} attend meeting at Germany month next\\
        \glt ‘I will (intend to) join a meeting in Germany next month.’ \\
        \hfill (Modern Thai, constructed)
    \ex \label{ex:5b} \textit{cà} future/promise/determination\\
        \textthai{ฉัน} \textthai{จะ} \textthai{ไม่} \textthai{ซื้อ} \textthai{ของ} \textthai{ร้าน} \textthai{นี้} \textthai{อีก} \textthai{ต่อไป}\\
        \gll \textit{chǎn} \textit{cà} \textit{mây} \textit{sʉ́ʉ} \textit{khɔ̌ɔŋ} \textit{ráan} \textit{níi} \textit{ìik} \textit{tɔ̀ɔpay}\\
        I \textsc{fut} not buy thing shop this again further\\
        \glt ‘I will (promise, am determined) not (to) buy any stuff again from this shop.’ \\
        \hfill (Modern Thai, constructed)
    \ex \label{ex:5c} \textit{cà} future/prediction\\
        \textthai{อากาศ} \textthai{ที่} \textthai{ปลอดโปร่ง} \textthai{จะ} \textthai{ทำให้} \textthai{เรา} \textthai{ฟื้นตัว} \textthai{จาก} \textthai{ไข้} \textthai{เร็ว} \textthai{ยิ่ง} \textthai{ขึ้น}\\
        \gll \textit{aakàat} \textit{thîi} \textit{plɔ̀ɔtpròoŋ} \textit{cà} \textit{thamhây} \textit{raw} \textit{fʉ́ʉntua} \textit{càak} \textit{khây} \textit{rew} \textit{yîŋ} \textit{khʉ̂n}\\
        air that clean \textsc{fut} make us recover from fever fast more up\\
        \glt ‘Clean air will (is expected to) help us recover from a fever faster.’ \\
        \hfill (Modern Thai, constructed)
    \ex \label{ex:5d} \textit{cà} future/irrealis\\
        \textthai{ถ้าผม} \textthai{จะ} \textthai{เปลี่ยนอาชีพ} \textthai{ผม} \textthai{จะ} \textthai{เป็น} \textthai{ครู}\\
        \gll \textit{thâa} \textit{phǒm} \textit{cà} \textit{plìan} \textit{aachîip} \textit{phǒm} \textit{cà} \textit{pen} \textit{khruu}\\
        If \textsc{1sg:m} \textsc{fut} change occupation \textsc{1sg.m} \textsc{fut} be teacher\\
        \glt ‘If I had changed careers, I would be a teacher (now).’ \\
        \hfill (Modern Thai, constructed)
    \z
\z

In \xref{ex:5}, \textit{cà} signals that the events described belong to a point in time in the future and at the same time implies certain concepts typically associated with the future tense, as indicated above. Evidently, in \xref{ex:5a}, the agent of the predicate with an action verb in the active voice (i.e., attend a meeting) is the sentential subject, which coincides with the speaker (i.e., ‘I’), an ideal condition for invoking the speaker’s intention. In \xref{ex:5b}, the semantic-syntactic features are similar to those of \xref{ex:5a}, with an added feature of negation of intention of any further patronage of the shop, which lends the interpretation of determination and promise. In \xref{ex:5c}, the sentential subject is non-sentient (i.e., air), thus not amenable to any interpretation involving conative modality such as intention, promise, determination, etc. The pragmatic inferences, based on the scientific knowledge of the causal relationship between health benefit and clean air, tend to suggest a prediction. Finally, as in \xref{ex:5d}, if the auxiliary \textit{cà} is used in the apodosis or protasis of a conditional sentence, especially of a hypothetical condition, it functions as an irrealis marker (cf. ‘hypothetical situation or putative action’ \cite{Noss1964}: 165; \cite{Soithurum2010}). This is in line with the observation that future is commonly associated with a prediction or extrapolation from the present state of the world (\cite{Dahl1985}: 103; \cite{comrie_identifying_1989}; \cite{Fleischman1982}: 24--31, 133--134; see also \cite{Kratzer1991}; \cite{Sarkar1998}; \cite{Sanso2006}: 243; \cite{Soithurum2010}; \cite{aikhenwald2010}: 81; \cite{StenzelGomezImbert2018}: 362, for individual instances). 

In addition to the usages of \textit{cà} exemplified above, \textit{cà} also signals diverse future-related modal notions, e.g., necessity, indeterminacy, possibility, certainty, proximation, constancy, etc., through univerbation with other lexemes, which contribute additional meanings of futurity. In this context, the univerbation phenomenon, frequently observed in Thai, warrants a brief exposition. When \textit{cà }forms a univerbated unit with another form (typically a modal), it is possible to treat the two modals in a complex form as operating separately, instead of treating the form as a single unit.\footnote{Our thanks go to the anonymous reviewer for raising this issue.} For instance, \textit{cà }and \textit{tɔ̂ŋ }(see \ref{ex:6a} below) each may be analyzed as signaling future and deontic necessity. This is a perfectly valid approach, especially considering that when \textit{cà }is omitted from \textit{cà-tɔ̂ŋ, }there may not be loss of deontic necessity (obligation) meaning, and conversely when \textit{tɔ̂ŋ }is omitted from \textit{cà-tɔ̂ŋ}, there may not be loss of futurity meaning. Such an analytic approach, however, does not capture the insights of the researchers who acknowledge the patterns of common co-occurrence of \textit{cà }and other modal elements, and therefore present them as more or less combined forms (even though \textit{cà }may be optional in certain cases). For instance, \citet[117]{Noss1964}, \citet[65--69]{Smyth2002}, \citet[130--134]{IwasakiIngkaphirom2005}, and \citet[58]{Soithurum2010}, among others, present certain combinations as unitized forms, either explicitly or implicitly, e.g., \textit{(cà) tɔ̂ŋ} ‘must’, \textit{khuan} \textit{(cà) }‘should/ought’, \textit{nâa} \textit{cà} ‘should/ought’, \textit{yàak} \textit{(cà)} ‘want to, would like to’, \textit{campen} \textit{cà }‘necessary to’, \textit{kamlaŋ cà} ‘be about to (happen)’, \textit{kʉ̀ap (cà)} ‘be about to (happen)’, \textit{thɛ̂ɛp (cà)} ‘be about to (happen)’, \textit{cuan (cà)} ‘be about to (happen)’, \textit{mák cà} ‘used to, would’, \textit{yɔ̂m cà} ‘likely’, \textit{chák cà} ‘beginning to’, \textit{nâa cà} ‘possibly’, \textit{khoŋ cà} ‘probably’, \textit{àat cà} ‘probably, perhaps’, etc. Therefore, while the treatment of such forms either analytically or holistically may be an unnecessarily granular issue, the fact that they tend to accompany other forms and that such a tendency is not applicable to all modal elements supports the practice of treating them as more or less unitized forms (note that the combined forms are represented with a hyphen between the two components). 

The usages of the \textit{cà}\textit{-}univerbated forms for signaling diverse future-related modal notions are exemplified in the following. The first set involves necessity as in \xref{ex:6}: 

\ea \label{ex:6} 
    \ea \label{ex:6a} \textit{cà-tɔ̂ŋ} future/necessity\\
        \textthai{คุณ} \textthai{จะต้อง} \textthai{จ่าย} \textthai{ค่าปรับ} \textthai{ที่} \textthai{สถานี} \textthai{ตำรวจ}\\
        \gll \textit{khun} \textit{cà-tɔ̂ŋ} \textit{càay} \textit{khâapràp} \textit{thîi} \textit{sathǎanii} \textit{tamrùat}\\
        you \textsc{fut} pay fine at station police\\
        \glt ‘You will (must, have to) pay the fine at the police station.’ \\
        \hfill (Modern Thai, constructed)

    \ex \label{ex:6b} \textit{khuan-cà} future/necessity\\
        \textthai{เธอ} \textthai{เปียก} \textthai{หมด} \textthai{เลย} \textthai{เธอ} \textthai{ควรจะ} \textthai{พก} \textthai{ร่ม} \textthai{มา} \textthai{ด้วย}\\
        \gll \textit{thəə} \textit{pìak} \textit{mòt} \textit{ləəy} \textit{thəə} \textit{khuan-cà} \textit{phók} \textit{rôm} \textit{maa} \textit{dûay}\\
        you wet all already you \textsc{fut} bring umbrella come together\\
        \glt ‘You are all wet. You should have brought an umbrella.’ \\
        \hfill (Modern Thai, constructed)
    \z
\z


Example \xref{ex:6a} above illustrates the usage of the future auxiliary \textit{cà-tɔ̂ŋ}, which also signals necessity. The additional morpheme \textit{tɔ̂ŋ} is itself a deontic auxiliary, a heteroseme of the verb \textit{tɔ̂ŋ} ‘to touch, be exposed to, agree’ (see Excerpt \ref{ex:3} of \textit{càk-tɔ̂ŋ} in Section \ref{sec:khammee:2.1} above). Since the phonologically longer form \textit{càk-tɔ̂ŋ} and the shorter form \textit{cà-tɔ̂ŋ} carry the identical function, it can be assumed that \textit{càk-tɔ̂ŋ} and \textit{cà-tɔ̂ŋ} are variants and that the latter is the more grammaticalized of the two (see Section \ref{sec:khammee:3.3.4} below). Similarly, \xref{ex:6b} exemplifies the use of \textit{khuan-cà}, an auxiliary for future and necessity, which denotes unfulfilled obligation. This compound incorporates the future \textit{cà} and the necessity modal auxiliary \textit{khuan}. Even though \textit{cà-tɔ̂ŋ} and \textit{khuan-cà} both denote future and necessity, the former is stronger in terms of the degree of necessity than the latter, analogous to the English deontic modals \textit{must} and \textit{should}. Another subtlety is that \textit{cà-tɔ̂ŋ} is more strongly future-oriented than \textit{khuan-cà}\textit{,} and thus in regret contexts as in \xref{ex:6b}, \textit{khuan-cà} is more appropriate. 

As shown above, the position of \textit{cà} in the combined forms varies (e.g., \textit{cà-tɔ̂ŋ} and \textit{khuan-cà}).\footnote{An anonymous reviewer raised this interesting issue, which may be related to scope differences.} It is generally accepted that the relative position of a form with reference to another typically carries the meaning-changing potential due to the scope effect, as in the famous contrast between \textit{ino-rssu-anguag} [person-big-little] ‘little giant’ and \textit{inu-ngua-rssuag} [person-little-big] ‘big midget’ in a dialect of Eskimo (\cite{SadockOlsen1976}, as cited in \cite{Bybee1985}: 96). \citet[]{Soithurum2010} identifies three different structural patterns of \textit{cà}: (i) before a verb, (ii) before or after an auxiliary, and (iii) after a complementizer (\textit{wâa} or \textit{thîi}), and briefly mentions, using the example of \textit{kamlaŋ cà, }the structural scope of \textit{kamlaŋ }and\textit{ cà }(\cite[72, 76]{Soithurum2010}). The combinations described in the present analysis all belong to (ii). While \citeauthor{Soithurum2010}'s classification system is useful, it is too imprecise to differentiate the types of, e.g., \textit{cà-tɔ̂ŋ} and \textit{khuan-cà}, and does not adequately explain scope differences. Considering that \textit{tɔ̂ŋ} and \textit{khuan}, for example, are both full-fledged auxiliaries of necessity, and that the relative position of \textit{cà }in \textit{cà-tɔ̂ŋ} and \textit{khuan-cà} cannot be reversed (*\textit{tɔ̂ŋ-cà} and *\textit{cà-khuan}), the fixed relative position of \textit{cà }is a great puzzle. It is possible that this frozen order represents the syntagmatic arrangement of the forms at the time of grammaticalization (cf. \cite{Givon1971}: 413). With the current paucity of relevant data, however, the issue requires further research.

The second set involves modals denoting variable degrees of certainty. 


\ea \label{ex:7}
    \ea \label{ex:7a} \textit{nâa-cà} future/probability (‘probably’)\\
        \textthai{ทุกอย่าง} \textthai{น่าจะ} \textthai{จบลง} \textthai{ด้วย} \textthai{ดี}\\
        \gll \textit{thúkyàaŋ} \textit{nâa-cà} \textit{còp.loŋ} \textit{dûay} \textit{dii}\\
        everything \textsc{fut} end by good\\
        \glt ‘Everything should end well.’ \\
        \hfill (\cite{Soithurum2010}: 79, modified)

    \ex \label{ex:7b} \textit{àat-cà} future/possibility (‘possibly’)\\
        \textthai{ครู} \textthai{อาจจะ} \textthai{สอน} \textthai{อยู่} \textthai{ที่} \textthai{ห้องเรียน} \textthai{สีแดง} \textthai{ตอนบ่าย} \textthai{พรุ่งนี้}\\
        \gll \textit{khruu} \textit{àat-cà} \textit{sɔ̌ɔn} \textit{yùu} \textit{thîi} \textit{hɔ̂ŋrian} \textit{sǐidɛɛŋ} \textit{tɔɔnbàay} \textit{phrûŋníi}\\
        teacher \textsc{fut} teach \textsc{prog} at classroom red.color afternoon tomorrow\\
        \glt ‘The teacher may teach in the red classroom tomorrow afternoon.’ \\
        \hfill (Modern Thai, constructed)

    \ex \label{ex:7c} \textit{khoŋ-cà} future/possibility (‘possibly’)\\
        \textthai{คุณตา} \textthai{คงจะ} \textthai{ปลูก} \textthai{ผัก} \textthai{อยู่} \textthai{หลังบ้าน} \textthai{เย็น} \textthai{นี้}\\
        \gll \textit{khuntaa} \textit{khoŋ-cà} \textit{plùuk} \textit{phàk} \textit{yùu} \textit{lǎŋbâan} \textit{yen} \textit{níi}\\
        grandpa \textsc{fut} plant vegetable at backyard evening this\\
        \glt ‘Grandpa may have planted vegetables at the backyard this evening.’ \\
        \hfill (Modern Thai, constructed)
    \z
\z



Example \xref{ex:7a} illustrates the use of \textit{nâa-cà} denoting future and probability. The added form \textit{nâa} is a modal adverb/auxiliary denoting possibility, which is inherent in future, and vice versa. The conceptual affinity may have motivated this univerbation process. As a future marker, \textit{nâa-cà} also indicates indeterminacy or uncertainty; thus, the described event is lower on the scale of certainty than \textit{cà-tɔ̂ŋ }derived from E-Necessity (high probability). The marker is also used in a regret situation, i.e., in describing an event that was once possible but was missed. Example \xref{ex:7b}, involving \textit{àat-cà}, in which \textit{àat }is also a modal adverb/auxiliary, indicates likelihood.\footnote{The lexeme \textit{àat} is also used as an adjective denoting ‘be brave, courageous, bold, be able’; it is likely, though not conclusively, a heteroseme of the modal adverb/auxiliary \textit{àat} ‘probable’.} Its degree of certainty is also lower than probability, as signaled by \textit{cà-tɔ̂ŋ}. Example \xref{ex:7c}, involving \textit{khoŋ-cà}, in which \textit{khoŋ} is a modal adverb/auxiliary of possibility, still bears a semantic connection with its adjectival heteroseme \textit{khoŋ} ‘be permanent, fixed’. The epistemic meaning (E-Possibility) of \textit{khoŋ-cà} assigns the word to be used as an evidential marker, as in \xref{ex:7c}. In this context, variability of the perceived degree of certainty as associated with the modals is evident among speakers of Thai, even conceding that the boundary between possibility and probability is often dependent on the speaker’s subjective judgment. For instance, \citet[137]{IwasakiIngkaphirom2005} state that \textit{tɔ̂ŋ }signals high confidence; \textit{nâa }and \textit{khuan }signal medium confidence (\textit{nâa }being more ‘subjective’ than \textit{khuan}); and \textit{khoŋ }and \textit{àat }signal low confidence. An anonymous reviewer stated that \textit{nâa,} \textit{àat }and \textit{khoŋ }are widely accepted to indicate a moderate degree of possibility similar to the English \textit{may} and \textit{might}. The judgment may be influenced by the propositional content as well. Minor perceptual disagreements notwithstanding, \textit{cà-}derived futurity markers may be said to signal E-Necessity and E-Possibility by virtue of involving epistemic modals in their combinations. 

\largerpage[1]

Another set of \textit{cà}-derived future modals consists of those marking ‘immediate future’ (\cite{Smyth2002}: 69; ‘imminent future’ \cite{HeineKutevaNarrog2017}: 10), as exemplified in the following: 

\ea \label{ex:8}
    \ea \label{ex:8a} \textit{kamlaŋ-cà} immediate future\\
        \textthai{เรา} \textthai{กำลังจะ} \textthai{กิน} \textthai{ข้าว}\\
        \gll \textit{raw} \textit{kamlaŋ-cà} \textit{kin} \textit{khâaw}\\
        we \textsc{fut} eat food\\
        \glt ‘We are/were about to eat.’ \\
        \hfill (\cite{Smyth2002}: 69, modified)
    
    \ex \label{ex:8b} \textit{cuan-cà} immediate future\\
        \textthai{เขา} \textthai{จวนจะ} \textthai{จม} \textthai{น้ำ} \textthai{แล้ว}\\
        \gll \textit{khǎw} \textit{cuan-cà} \textit{com} \textit{náam} \textit{lɛ́ɛw}\\
        he \textsc{fut} drown water \textsc{pfv}\\
        \glt ‘He is about to drown.’ \\
        \hfill (Modern Thai, constructed)
    
    \ex \label{ex:8c} \textit{kʉ̀ap-cà} immediate future\\
        \textthai{ฉัน} \textthai{ทำ} \textthai{การบ้าน} \textthai{เกือบจะ} \textthai{เสร็จ} \textthai{แล้ว}\\
        \gll \textit{chǎn} \textit{tham} \textit{kaanbâan} \textit{kʉ̀ap-cà} \textit{sèt} \textit{lɛ́ɛw}\\
        I do homework \textsc{fut} finish \textsc{pfv}\\
        \glt ‘I will soon finish (almost finished) my homework.’ \\
        \hfill (Modern Thai, constructed)
    
    \ex \label{ex:8d} \textit{thɛ̂ɛp-cà} immediate future\\
        \textthai{ฉัน} \textthai{หิว} \textthai{แทบจะ} \textthai{ตาย}\\
        \gll \textit{chǎn} \textit{hǐw} \textit{thɛ̂ɛp-cà} \textit{taay}\\
        I be.hungry \textsc{fut} die\\
        \glt ‘I’m so hungry (I feel like) I’m going to die.’ \\
        \hfill (\cite{HigbieThinsan2002}: 90, modified)
    
    \ex \label{ex:8e} \textit{klây-cà} immediate future\\
        \textthai{งาน} \textthai{นี้} \textthai{ใกล้จะ} \textthai{เสร็จ} \textthai{แล้ว}\\
        \gll \textit{ŋaan} \textit{níi} \textit{klây-cà} \textit{sèt} \textit{lɛ́ɛw}\\
        work this \textsc{fut} finish \textsc{pfv}\\
        \glt ‘The job is nearly done.’ \\
        \hfill (\cite{HigbieThinsan2002}: 89, modified)
    \z
\z

The examples in \xref{ex:8} involve future markers that also denote that the event being described will occur in the very near future. However, as \citet[69]{Smyth2002} noted with \xref{ex:8a}, the use of the future marker is not restricted to referring to a future time in relation to the speech time, but can also be used to refer to a later time in relation to the time of reference in the past. In the literature, the notion of immediacy is often separately labeled as ‘proximative’ or ‘prospective’ (see \cite{Heine1994}; \cite{Johanson2017}; \cite{Nevskaya2005}, \citeyear{Nevskaya2015}; \cite{Kuteva2019}, for terminological issues).\footnote{\citet[]{Johanson2017} distinguishes prospective and proximative: prospective being an aspect marker signaling a perspective on events foreseen to occur subsequent to some reference time, and proximative being an ‘actionality’ marker referring to a pre-phase (preliminary or preparatory) of a possible event conceived of as imminent or impending, i.e., as ‘about to occur’ (as cited in \cite{Nevskaya2015}). For similar issues, see \citet[64]{Comrie1976}; \citet[111--112]{Dahl1985}; among others.} Example \xref{ex:8a} involves the compound modal \textit{kamlaŋ-cà}, in which \textit{kamlaŋ} is a progressive marker (note that \citealt[149]{IwasakiIngkaphirom2005} label \textit{kamlaŋ-cà} as an ‘inceptive’ (i.e., ‘is about to do’) aspect marker). Example \xref{ex:8b} illustrates the use of the immediate future marker \textit{cuan-cà}, which involves the adverb of proximity \textit{cuan} ‘almost, nearly, close to’ (note that \citet[149]{IwasakiIngkaphirom2005} label \textit{cuan-cà} as a ‘pre-inceptive’ (i.e., ‘is almost done’) aspect marker). The closeness in degree or space is used to denote temporal proximity. Similarly, \xref{ex:8c}, involving the immediate future \textit{kʉ̀ap-cà}, makes use of the degree adverb/auxiliary \textit{kʉ̀ap} ‘almost, nearly, close to’. Likewise, \xref{ex:8d} makes use of the degree adverb \textit{thɛ̂ɛp} ‘almost, nearly, barely’, and \xref{ex:8e} makes use of \textit{klây} ‘(be) near, beside’, a polyfunctional exponent in the adjective, adverb, and preposition word class. 

Just as \textit{cuan}, \textit{kʉ̀ap}\textit{, }\textit{thɛ̂ɛp}\textit{, }and \textit{klây} are synonymous, so are the derived future modals \textit{cuan-cà}, \textit{kʉ̀ap-cà}\textit{, }\textit{thɛ̂ɛp-cà}\textit{, }and \textit{klây-cà}. There is no perceived difference in the degree of proximity in these immediate future markers. However, as compared to \textit{kamlaŋ-cà} in \xref{ex:8a}, the degree of immediacy signaled by these four, \xref{ex:8b} through \xref{ex:8e}, is greater. 

There is another, sui generis future marker, i.e., \textit{mák-cà}, the usage of which is exemplified below: 

\ea \label{ex:9}
    \textit{mák-cà }future/habitual\\
    \textthai{น้ำหวาน} \textthai{มักจะ} \textthai{ทาน} \textthai{ขนม} \textthai{ใน} \textthai{รถ}\\
    \gll \textit{Namwan} \textit{mák-cà} \textit{thaan} \textit{khanǒm} \textit{nay} \textit{rót}\\
    [name] \textsc{fut} eat snack in car\\
    \glt ‘Namwan would eat (always eats) snacks in the car.’ \\
    \hfill (Modern Thai, constructed)
\z

The future marker \textit{mák-cà} is a compound modal involving the degree adverb \textit{mák} ‘frequently, often’. By virtue of the involvement of the full-fledged future marker \textit{cà}, the compound modal \textit{mák-cà} marks the future, and by virtue of the frequency adverb \textit{mák}, it marks the high frequency of repeated occurrence, hence a habitual future. Thus, the event being described is thought to have occurred habitually in the past, at the present time, or in the future. In Modern Thai, the semantic bleaching occurred to such an extent that \textit{mák-cà} is also used for habitual occurrence of past events (as in the English modals \textit{used to }or \textit{would}). This seems to be due to the fact that the speaker’s assessment that some events will occur in the future with regularity is often based on the same events having happened with regularity in the past and in the present. If that is the case, using a future habitual marker for past habitual events is an instance of reasonable extension along the timeline (cf. the English habitual \textit{will} and \textit{would }based on prediction, \cite{QuirkEtAl1985}: 228--229). 

\largerpage[0.5]

\subsection{Adverbials}\label{sec:khammee:2.3}

There are other, less productive, less entrenched, future markers (\cite{Supanvanich1973}), i.e., immediate future markers originating from ‘almost, close, nearly’, i.e., \textit{cuan}\textit{, }\textit{kʉ̀ap}\textit{, }\textit{thɛ̂ɛp}\textit{, }and \textit{klây}\textit{, }as exemplified below:

\ea \label{ex:10}
    \ea \label{ex:10a} Immediate future\\
    \textthai{ฝน จวน ตก แล้ว กลับบ้าน กัน}\\
    \gll \textit{fǒn} \textit{cuan} \textit{tòk} \textit{lɛ́ɛw} \textit{klàp} \textit{bâan} \textit{kan}\\
    rain \textsc{fut} fall \textsc{pfv} return home \textsc{ptcl.hort}\\
    \glt ‘It’s going to (almost started to) rain. Let’s go back home.’ \\
    \hfill (Modern Thai, constructed)\\

    \ex \label{ex:10b}
    \textthai{เขา เกือบ}/\textthai{แทบ}/\textthai{ใกล้} \textthai{จม น้ำ แล้ว}\\
    \gll \textit{khǎw} \textit{kʉ̀ap/thɛ̂ɛp/klây} \textit{com} \textit{náam} \textit{lɛ́ɛw}\\
    he \textsc{fut} drown water \textsc{pfv}\\
    \glt ‘He is about to drown.’ \\
    \hfill (Modern Thai, constructed)
    \z
\z


The examples in \xref{ex:10} illustrate \textit{cuan}, \textit{kʉ̀ap}\textit{,} \textit{thɛ̂ɛp}\textit{,} and \textit{klây} functioning as future markers. These four lexemes are those that formed compound modals for the immediate future (see \ref{ex:8b} through \ref{ex:8e} above). In other words, these examples differ from the previously addressed \textit{cuan-cà}\textit{,} \textit{kʉ̀ap-cà}\textit{, }\textit{thɛ̂ɛp-cà} and \textit{klây-cà} only in that they do not involve the primary future marker \textit{cà}; the futurity is marked solely by the proximity markers denoting ‘almost, nearly, close’. 

One peculiarity of these adverb-based immediate future markers is that they often, but not necessarily, co-occur with \textit{lɛ́ɛw}\textit{,} the perfective aspect particle, also functioning as a sentence ending particle of diverse discourse functions (note that \ref{ex:10a} is bi-clausal (or bi-sentential), the first one ending with \textit{lɛ́ɛw}), which originated from the verb \textit{lɛ́ɛw} ‘finish, be done’ (\cite{IwasakiIngkaphirom2005}: 149).\footnote{The immediacy effect created by \textit{lɛ́ɛw} is evident in the fact that when the primary future marker \textit{cà} occurs with \textit{lɛ́ɛw}, the sentence means that the event being described will occur very shortly.} The combination of the perfective \textit{lɛ́ɛw} with the proximity degree markers signals that something will not only occur soon but will also be completed soon, and thus creates the sense of definite and complete occurrence of the event within a short period of time. 

Treating temporal adverbials as future markers, as \citet{Supanvanich1973} does, raises an important issue that warrants mention, as an anonymous reviewer pointed out. According to the reviewer, adverbials modify actions with respect to their degree of progress, but the described actions may not be actualized. As researchers have observed, futurity is never a purely temporal concept but it necessarily includes an element of prediction or some related modal notion (\cite{Lyons1977}: 677; see also \cite{Salkie2010}: 1787). There is always an inherent uncertainty about whether an event will take place or not (\cite{Dahl1985}, \citeyear{Dahl2000}; \cite{JaeggiEtAl2022}). Thus, a future marker does not guarantee the occurrence of a described event or the presence of a described state of affairs in the future; it only signals that an event or state is, unless stated otherwise, to come about at a time later than the speech time or referenced time. The detachment of a future-marked event from its actualization in reality is inherent in the notion of future, and this trait closely links the future marker to irrealis (\cite{Noss1964}; \cite{Soithurum2010}: 110). \citet[165]{Noss1964}, in his exposition of \textit{cà-,} states that the meaning of \textit{cà- }is ‘hypothetical situation or putative action.’ In this regard, the analysis by \citet[123]{IwasakiIngkaphirom2005} of \textit{cà }as a marker of ‘challengeability’ (\cite{Givon1982}) is particularly noteworthy. In other words, a future event is assumed to be challengeable by virtue of not having been realized and thus may be marked by \textit{cà. }Such an analysis can shed light on the usage of \textit{cà }for non-temporal marking, the discussion of which is beyond the immediate scope of the present paper (see \cite{IwasakiIngkaphirom2005}: 126--127 for more exposition).

The issue of whether or not the forms marking the degree of progress may be regarded as future markers is debatable. Two points are relevant to this issue: (i) the Mandarin Chinese temporal lexeme \textit{yào} ‘immediately’ (as well as ‘want’) functions as a ‘lexical futurity’ marker (\cite{LiThompson1989}: 175), and (ii) the general notion of futurity may include temporal adverbials, which are lexical, as well as proximative and prospective adverbials, which are often on the boundary of lexis and grammar. In this light, the adverbials that denote temporal proximity addressed here may well be regarded, at least, as ‘peripheral’ markers of futurity. 



\subsection{Particles}\label{sec:khammee:2.4}

\largerpage[1]

Other linguistic forms have additionally been claimed to be future markers in Thai. For instance, \citet[]{Supanvanich1973} asserts that certain sentence ending particles are future markers, e.g., \textit{lá}\textit{, }\textit{ná}\textit{, }and \textit{lá-ná}. These future markers are subject to controversy and, even if qualified as such, are undoubtedly marginal forms. Such uses are exemplified in the following:

% \newpage

\ea \label{ex:11}
    Immediate future\\
    \textthai{ฉัน นอน ละ} / \textthai{นะ} / \textthai{ละนะ}\\
    \gll \textit{chǎn} \textit{nɔɔn} \textit{lá / ná / lá-ná}\\
    I sleep \textsc{ptcl}\\
    \glt ‘I (will) sleep now.’ \\
    \hfill (\cite{Supanvanich1973})
\z

The example in \xref{ex:11} involves the sentence ending particles \textit{lá, ná, }and \textit{lá-ná.} The primary function of these particles, when used in any sentence, is to draw the addressee's attention to the information being provided and indicate this information as noteworthy. The notion of futurity associated with these forms seems to be only pragmatically inferred, i.e., not yet grammaticalized. In other words, when the example sentence \xref{ex:11} is spoken, and the speaker is not yet in bed, the addressee, noticing the discrepancy, reconciles it by inferring that the proposition (\textit{C}\textit{hǎn nɔɔn} ‘I sleep’), presented as a piece of ‘noteworthy’ information (as indicated by \textit{lá/ná/lá-ná})\textit{,} is describing a situation ‘as good as the present’, hence the immediate future. Put another way, in \xref{ex:11} above, there exists a gap between the speaker's alleged claim (i.e., ‘I am sleeping’) and the hearer's perceived reality (i.e., ‘You are not sleeping now’), which triggers the pragmatic inference (‘You will go to sleep immediately’). The motivation behind describing a future event as one that is presently occurring is conceptually straightforward, and is classifiable as an instance of hyperbole (cf. ‘That does it! I’m leaving’ \cite{QuirkEtAl1985}: 215).\footnote{The role of hyperbole in grammaticalization has been reported in a series of works by \citeauthor{Rhee2003} (\citeyear{Rhee2003}, \citeyear{Rhee2009}, \citeyear{Rhee2022}) on the Korean language.} This pragmatic inference seems to have triggered the functional reanalysis of the sentence ending particles as immediate future markers, or at least as forms with a future connotation. It is to be noted that this inference of future time reference is triggered by only these sentence ending particles out of numerous particles in Modern Thai. Therefore, this is different from the use of the simple present tense for future, which is widely attested across languages (cf. \cite{QuirkEtAl1985}: 182--183 for English). Incidentally, the variations by means of the three sentence ending particles do not create any significant differences in meaning, but most Thai speakers understand that when the longer particle (\textit{lá-ná}) is used instead of the shorter ones (\textit{lá} and \textit{ná}), the proposition carries stronger illocutionary force. 

\largerpage[1]

\section{Discussion} \label{sec:khammee:3}

In the preceding section, we have seen how Thai future markers are used. Based on the description, we will discuss a few issues involved in their development that bear theoretical significance, e.g., morphosyntactic motivation, conceptual motivation, and grammaticalization parameters.


\subsection{Morphosyntactic Motivation}\label{sec:khammee:3.1}

The morphosyntactic motivation behind the development of Thai futures comprises two primary factors, one typological, and the other idiosyncratic. The typological factor relates to pragmatic orientation and verb serialization, and the idiosyncratic factor relates to the preference of creating polylexemic units in lexicalization and grammaticalization. 

Most functionalists hold that languages differ in terms of the degree to which they utilize pragmatic factors in their linguistic representation. For instance, \citet[276]{Huang2000} asserts that ‘Some languages are more pragmatic than others’ (for similar positions, see \cite{LiThompson1976}, \citeyear{LiThompson1979}; and on Thai, see \cite{bamgroongraks1987}; see also \cite{OGrady2005}, \citeyear{OGrady2008}, for the ‘emergentist’ approach to linguistic structure). Thai lacks many grammatical forms that are found in many languages, notably linking words. Thus, the semantic and syntactic relationships between constituents are often identified from the linguistic and situational context. This pragmatic orientation is related to the proliferation of verb serialization, in which verbs are strung together asyndetically. \citet[160]{Diller2006}, while illustrating an overload of verb-based detail in informal Thai, presents an example of a seemingly extreme case: ‘Go. Hurry exit go buy return come finish peel pass’ (for ‘Then hurry off and buy some. When you get back, peel them – get on with it’). The serial verb construction is so susceptible to grammaticalization that it is often regarded as the ‘seed’ of grammaticalization (cf. \cite[15--16]{DeLancey1991}). Indeed, all serial verb languages seem to develop exponents of grammatical functions out of serial verb constructions by way of reanalysis (\cite{Bynon1985}), or of pragmatics, i.e., interaction between conceptualization of world events and iconicity in the language (cf. \cite{Bruce1988}; \cite{DeLancey1991}; \cite{Lord1995}; \cite{Rhee1996}; \cite{aikhenwald_dixon2006} and works therein). With the prevalence of verb serialization, Thai exhibits numerous instances of grammaticalization from serialized verbs (\cite{Vis1978}; \cite{Thepkanjana1986}; \cite{Diller2006}; \cite{Park2011}). 

The development of the primary future tense marker \textit{cà} (and its older \textit{càk}), which originated from a verb, seems to have occurred in verb serialization. The likely lexical origin of the marker is the Sino-Thai lexeme \textit{càk} ‘to know, recognize’, which survives in contemporary Thai as \textit{rúucàk} ‘to know a person, place, etc.’ (see Section \ref{sec:khammee:2.1} above and Section \ref{sec:khammee:3.2} below). Thus, the source construction of the future marker is, for example, [I + know + return + home], from which the reinterpreted construction [I + will + return + home] emerged (Note that Thai is a head-initial language and the auxiliary precedes the verb).

Secondly, Thai has a strong idiosyncratic preference for the creation of polylexemic or polysyllabic units in lexicalization and grammaticalization (Note that native Thai lexemes are predominantly monosyllabic, \cite{Smyth2002}: 1). In other words, Thai often resorts to recycling a single lexeme to create similar or slightly modulated meanings by reduplicating the base form, reinforcing form and meaning by means of the juxtaposing of synonyms, and/or recruiting additional forms, such as particles, adverbs, etc. (cf. ‘polysemy strategy’ \cite{KhammeeRhee2022}; \cite{RheeKhammee2022}; \cite{ThirakunkovitRhee2023}). This preference of ‘addition’ is also applicable to the sentence level; various sentence ending particles are (often multiply) added to the end of the sentence, with the primary function being to soften the illocution and make the assertion or statement more polite. The speech style of ending a sentence without a sentence-final particle is stigmatized as \textit{mây-mii-hǎaŋ-sǐaŋ} ‘having no tail of sound’ (\cite{khammee2024demonstratives}). 

The preference for creating polylexemic or polysyllabic units by means of the above-mentioned strategies is responsible for the creation of several \textit{cà-}derived future markers, e.g., \textit{cà-tɔ̂ŋ, khuan-cà, nâa-cà, àat-cà, khoŋ-cà, kamlaŋ-cà, cuan-cà, kʉ̀ap-cà, thɛ̂ɛp-cà, klây-cà, mák-cà,} etc. (see Section \ref{sec:khammee:2.2} above). When multiple modals occur in a single sentence, they tend to occur next to each other in the pre-verbal slot, and thus their positional closeness, as well as the general tendency to create polylexemic, polysyllabic units, is likely to increase the bond between them (hence, \textit{cà-tɔ̂ŋ} ‘will’ (< ‘will-must’)). Similarly, adverbials often occur next to the future marker. In general, there is a small group of adverbials that occur sentence-finally (but before sentence ending particles), e.g., \textit{samə̌ə} ‘always’, \textit{bɔ̀y} ‘often’, and \textit{rʉ̂ay} ‘all the time’ (\cite{IwasakiIngkaphirom2005}: 103), but others are more positionally flexible and often occur at the pre-auxiliary position. Hence, the habitual future \textit{mák-cà} ‘will’ (< ‘frequently-will’), the immediate future \textit{cuan-cà} ‘be about to, will’ (< ‘almost-will’), etc.

Apart from the two primary factors elaborated above, a mention is in order with respect to the immediate future markers that are heterosemous with degree adverbs, i.e., \textit{cuan} ‘frequently’,\textit{ kʉ̀ap, thɛ̂ɛp}, and \textit{klây} (see Section \ref{sec:khammee:2.3} above). These future markers are less established, as evident in the fact that few researchers include them in the inventory of future forms (e.g., \cite{Supanvanich1973}). The syntagmatic position of degree adverbials is flexible, but the pre-verbal position is favored, because they are modifiers of the verb or verb phrase. Thus, in this configuration, degree adverbials occupy the same position as the auxiliaries, including future markers such as \textit{cà}, which, coupled with conceptual motivation (see below), seems to have triggered the functional reinterpretation and structural reanalysis of the degree adverbials into future auxiliaries.\footnote{See Section \ref{sec:khammee:2.2}, however, for unsolved morphosyntactic positionality issues.}



\subsection{Conceptual Motivation}\label{sec:khammee:3.2}

Semantic change in grammaticalization is not random but, rather, is motivated (‘the source determination hypothesis’ and ‘the universal path hypothesis’; \cite{BybeePerkinsPagliuca1994}). Conceptual motivation and the gradient nature of linguistic change lead to the formation of a conceptual network of lexemes and grammatical forms (\cite{HeineClaudiHunnemeyer1991}; \cite{Nuyts1992}; ‘semantic map’ \cite{vanderAuwera1998}; \cite{Croft2001}; \cite{Narrog2009}). 

We have noted that the source lexeme of the Thai primary future marker \textit{cà} is likely the Sino-Thai verb \textit{càk} ‘to know, recognize’ (see Section \ref{sec:khammee:2.1} above) and that around the time when it became grammaticalized into a future auxiliary in the Sukhothai period (1238--1347 CE), its meaning became ‘to intend’ (see Section \ref{sec:khammee:2.2} above). Crosslinguistic research shows that the future develops from diverse sources, e.g., Change-of-state, ‘come to’, ‘go to’, ‘love’, B-Necessity, D-Necessity, H-Possessive, ‘take’, ‘then’, ‘tomorrow’, Venitive, and ‘want’ (\cite{Kuteva2019}). Since B-Necessity denotes ‘intention’, the lexical source of the Thai path conforms to the crosslinguistic pattern (also the path proposed by \cite{BybeePerkinsPagliuca1994}). Furthermore, the lexical source ‘know’ is known to develop into PI-Possibility (Ability) and further to Habitual or Permission, or to Evidential or Experiential (\cite{Kuteva2019}). Mandarin Chinese has no grammatical tense, but the verb \textit{hui} ‘know’ can indicate futurity (\cite{BybeePerkinsPagliuca1994}: 265; \cite{LiThompson1989}: 183). Similar to this grammaticalization scenario, the Medieval Chinese verb \textit{jie} ‘know, know of’ developed into the auxiliary of PI-Possibility (Ability), which describes a notion closely related to the future (\cite{LongEtAl2012}). The Thai future marker \textit{cà} marks Possibility (E-Possibility), closely related to PI-Possibility (Ability), as well as Necessity (i.e., E-Necessity, for Probability; D-Necessity, for Obligation; and B-Necessity, for Intention), often with the help of co-occurring forms. The states of affairs in Thai fit with the observed crosslinguistic patterns.

It is noteworthy that once a lexical form or a construction begins to mark future tense, the marker, in most cases, seems to further acquire diverse modal functions, such as Necessity, Certainty, Probability, etc., even though the focus may be selective and different. For instance, the Thai future acquired the modal function of Indeterminacy, which is likely to have derived from the characteristic of the future as inherently indeterminate and uncertain. 

The adverbials of proximity denoting ‘almost, nearly’, either alone or with the primary future \textit{cà}, have developed into future markers (see Section \ref{sec:khammee:2.2} and Section \ref{sec:khammee:2.3} above). The development must have been enabled by the general proximity (typically spatial and degree) metaphorically projected to temporal proximity. The original semantics still persists in the grammaticalized markers in that these are immediate future markers (rather than remote future or general future) and that their functions are largely restricted to marking immediate future, not extended to other modal meanings.

We also noted that some sentence ending particles have developed into future markers, albeit marginally (see Section \ref{sec:khammee:2.4} above). Considering that the primary function of these markers is to signal that the speaker is ‘informing’ the addressee by providing noteworthy information, we argued that the addressee’s attempt to fill the gap between the information and reality is responsible for the emergence of the function of marking immediate future. 

All these conceptual considerations and their interrelatedness may be sketched, as in Figure \ref{fig:1}.

\begin{figure}
    \centering
    \includegraphics[width=0.8\linewidth]{figures/khammee/figure1.pdf}
    \caption{Conceptual network of Thai future markers}
    \label{fig:1}
\end{figure}

\subsection{Grammaticalization Parameters}\label{sec:khammee:3.3}

A brief mention to Thai documentation is in order before we discuss diachronic development. The documentation of the Thai language has a rather shallow historical depth. \citet[]{Haarmann2012}, \citet{Leppert1992}, and others state that it is possible that the early literature of the Thai people may have been written in Chinese, but no historical record of the Siamese thus far refers to this earlier literature. The Thai alphabet was established in the 13\textsuperscript{th} century CE, and the earliest literature dates back to the Sukhothai period (13\textsuperscript{th} to mid--14\textsuperscript{th} century), chiefly surviving in the form of stone inscriptions (e.g., the King Ramkhamhaeng inscription of 1292). Classical literature dating from the Ayutthaya period (1351--1767) is written in verse but most of it was lost in the 18\textsuperscript{th} century in the military conflict with Burma (Myanmar). Thus, the extant literature is recent, mostly dating back to the 19\textsuperscript{th} century. One obvious implication of this shallow historical depth in grammaticalization studies is that the developmental trajectories of grammatical forms cannot be easily or conclusively established. Thus, following the tradition of researching languages with little or no historical documentation (\cite{TraugottHeine1991}: 6), we can project diachronic grammaticalization to synchronic data and, conversely, synchronic variations to the reconstruction diachronic developmental paths.

The relative paucity of historical data notwithstanding, we will review the grammaticalization scenarios from the available resources with respect to the four grammaticalization parameters suggested in \citet[]{Kuteva2019}, i.e., desemanticization, extension, decategorialization, and erosion. 


\subsubsection{Desemanticization}\label{sec:khammee:3.3.1}

Also commonly called ‘semantic bleaching’ in figurative terms, desemanticization refers to the loss in meaning content. As briefly alluded to in the preceding exposition, the source lexeme \textit{càk }has undergone desemanticization to such an extent that most speakers of contemporary Thai are unable to recognize the connection between the verb \textit{càk }and the future marker \textit{cà}. The implication of cognizance, i.e., ‘know’, is nearly completely lost and its trace is only visible in the verb \textit{rúucàk} ‘to know a person, place, etc.’ Even as a future marker, \textit{cà} does not make a substantial semantic contribution to the \textit{cà}-compounds. For instance, the semantic difference between the compound and non-compound is often imperceptible (e.g., \textit{àat} ‘might’ and \textit{àat-cà} ‘might’; \textit{khoŋ} ‘possibly’ and \textit{khoŋ-cà} ‘would’; \textit{nâa} ‘probably’ and \textit{nâa-cà}‘likely’; \textit{khua}n ‘should’ and\textit{ khuan-cà }‘should’; \textit{kʉ̀ap} ‘almost’ and \textit{kʉ̀ap-cà} ‘almost’; etc.). 

On the other hand, the meanings of general degree and spatial proximity as associated with the proximity adverbs denoting ‘almost’ have been bleached, but the bleaching is not substantial because the immediate future is still based on the notion of temporal proximity. The sentence ending particles, by virtue of having no ‘semantic’ content, cannot be discussed in this regard. 

\subsubsection{Extension}\label{sec:khammee:3.3.2}

Also known as context generalization, extension refers to the use of a form in new contexts. With respect to the future \textit{cà}, the source lexeme was a verb denoting the presence of knowledge, which later denoted intention, and after grammaticalization the form has become an auxiliary of the future. Through grammaticalization, the usage context has been extended because there is no more semantic constraint and the form can modify a proposition (i.e., a clause). 

Similarly, the degree adverbs that, before grammaticalization, could be used as modifiers of verbal, adjectival, and predicate constituents, can now modify propositions, a clear instance of usage context extension. 

The sentence ending particles that could be used as modifiers of a clause and function as a discourse operator (as markers of the speaker’s diverse stances) can now be used as modifiers of a clause with respect to the temporal dimension of the proposition, which also constitutes an instance of usage context extension. 

\subsubsection{Decategorialization}\label{sec:khammee:3.3.3}

Decategorialization refers to the loss in morphosyntactic properties characteristic of lexical or other less grammaticalized forms. Decategorialization typically involves, but is not restricted to, the primary categories. We will look at the source verb \textit{càk} and the adverbs\textit{ cuan, kʉ̀ap, thɛ̂ɛp}, and \textit{klây}.

The verb \textit{càk}, by virtue of being a verb, could carry diverse verbal features, such as degree, quality, frequency, and manner modification, e.g., ‘fast’, ‘well’, ‘everyday’, ‘diligently’, ‘thoroughly’, etc. However, after its grammaticalization into a future marker, the form cannot be marked by such modifiers. The ability to c-select a theme as its complement, a feature of a transitive verb, has also been lost.

Similarly, the proximity adverbials, before their grammaticalization into future markers, could function as modifiers of verbs, adjectives, and predicates. After grammaticalization, such adverbial features are lost and can function only as modifiers of a clause. Thus, the grammaticalization of future markers in Thai clearly shows decategorialization. 

\subsubsection{Erosion}\label{sec:khammee:3.3.4}

Also known as phonetic reduction, erosion refers to the loss in phonetic substance. Even though most grammaticalizing forms across languages undergo phonological reduction, often by virtue of increased frequency, Thai future markers are not eroded to a significant level, a characteristic of an isolating language (cf. \cite{Matisoff1991}; \cite{bisang1996}, \citeyear{Bisang2011}; \cite{Enfield2009}; \cite{ansaldo2018}). For instance, the verb \textit{càk} did not lose any substance in its development into the future marker \textit{càk}. As it developed into the future marker \textit{cà}, some minimal loss has occurred. 

Similarly, the proximity adverbials have not undergone phonological erosion either, in the course of their development into future markers. Their syntagmatic properties as well as their phonological shapes (including tonal features) are retained even after the grammaticalization. The only change that occurred is their functional reinterpretation. The same is applicable to sentence ending particles: they retain all syntagmatic and phonological properties after their grammaticalization. The grammaticalization only lies in the reinterpretation of a discourse-pragmatic marker into a tense-marking modal auxiliary.

There is one noteworthy aspect in the grammaticalization of future markers in Thai, which may be extendable to Thai grammaticalization in general. It relates to the reverse of erosion, i.e., reinforcement. As indicated in Section \ref{sec:khammee:3.1}, Thai has a strong tendency to reinforce a linguistic (particularly grammatical) form. The common strategies recruited in the realization of this inclination include reduplicating the base form, reinforcing form and meaning by means of juxtaposing synonyms, and recruiting additional forms, such as particles, adverbs, etc. (see Section \ref{sec:khammee:3.1} above). For instance, in their research on Thai allatives numbering 26, \citet[]{RheeKhammee2022} report that there are many forms that recycle the base form by reinforcing it with semantically similar forms, e.g., \textit{hây} ‘to/for’, \textit{hây-kɛ̀ɛ} ‘to/for’, \textit{hây-kàp} ‘to/for’ all involving \textit{hây} ‘to/for’; \textit{con} ‘to’, \textit{con-thʉ̌ŋ} ‘to’, \textit{cùap-con} ‘to’, \textit{con-kràthâŋ-thʉ̌ŋ} ‘to’, \textit{tràap-con} ‘to’, \textit{tràap-con-thʉ̌ŋ} ‘to’, all involving \textit{con }‘to’, etc. (see also \cite{khammee2025goal} for a similar analysis). Similarly, \citet[]{ThirakunkovitRhee2023}, in their research on Thai intensifiers, note that out of 23 intensifiers as many as 14 of them involve \textit{mâak} ‘very, very much’, the primary intensifier in contemporary Thai, e.g., \textit{mâak} ‘very’, \textit{mâak-maay} ‘very’, \textit{mâak-maay-kàay-kɔɔŋ} ‘very’, \textit{mâak-kəən-pay} ‘very, too much’, \textit{mâak-con-kəən-pay} ‘very, too much’, etc. These forms, though composite in makeup, are perceived as ‘univerbated’ units (\cite{Lehmann2015}) by Thai speakers. Even though further research is required to reach a solid generalization, Thai seems to favor semantic and phonological reinforcement to increase the salience, especially as a grammaticalizing form becomes semantically weak or loses novelty through frequent use.


\section{Summary and Conclusion}\label{sec:khammee:4}

Contemporary Thai has a number of forms that mark futurity. In addition to the primary marker \textit{cà} and its older form \textit{càk}, there are secondary markers of futurity, e.g., proximity adverbials denoting ‘almost’ and sentence ending particles signaling informative noteworthiness, which are used either alone or in combination with other forms. An investigation of their morphosyntactic development reveals some influence of the typological and idiosyncratic features of Thai, such as a high level of pragmatic orientation, extensive verb serialization, and preference for polylexemic units in lexicalization and grammaticalization, among others. Their semantic developments are also well motivated in terms of conceptual relatedness between the source forms and the grammaticalized forms. A review of grammaticalization scenarios with reference to grammaticalization parameters, e.g., desemanticization, extension, decategorialization, and erosion, shows that the Thai future markers have been grammaticalized to varying degrees. One peculiarity of Thai future grammaticalization is that Thai favors semantic and phonological reinforcement to increase the salience, especially as a grammaticalizing form becomes semantically weak or loses novelty through frequent use, thus countering semantic bleaching and formal erosion.





\section*{Abbreviations}
\begin{tabularx}{.5\textwidth}{@{}lQ@{}}
\textsc{fut} & future \\
\textsc{hon} & honorific \\
\textsc{hort} & hortative \\
\textsc{m} & masculine \\
\textsc{nmlz} & nominalizer \\
\end{tabularx}%
\begin{tabularx}{.5\textwidth}{@{}lQ@{}}
\textsc{pfv} & perfective \\
\textsc{prs} & present \\
\textsc{prog} & progressive \\
\textsc{ptcl} & particle \\
\textsc{sg} & singular \\
\end{tabularx}

\section*{Acknowledgements}
Special thanks go to Prof. Dr. Stefan Hartmann and Lena Schnee for their kind support, to Prof. Dr. Seongha Rhee for his guidance and support, to an anonymous reviewer for meticulously checking the content and providing suggestions, and to Prof. Edgar C. Gordyn for reading an earlier version and helping improve it in style and content. This research was supported in part by the research fund of University of Phayao.

\sloppy
\printbibliography[heading=subbibliography,notkeyword=this]
\end{document}



